\section{\gai Bottlenecks}
\label{section:33-literature-gai-bottlenecks}

A \textit{bottleneck}, or \textit{system constraint} is any resource whose capacity is \nrscomment{less than or equal to} equal to or less than the demand placed on it~\cite{Goldratt2004}. 
Understanding the common bottlenecks in \gai systems helps design the tasks and datasets that will \nrscomment{discriminate \gai system...} be discriminative in judging \gai system performance. 
This section briefly reviews the prior work in identifying the bottlenecks known to occur \nrscomment{in} during \gai, 
\nrscomment{"which inform our benchmark design methodology" or something like that. skip the bit about elevate and exploit}
setting the scene for the discussion of acceleration techniques to \textit{exploit} and \textit{elevate} those constraints in section~\ref{section:34-literature-gai-acceleration}.

\nrscomment{this first sentence doesnt really fit, just cut it}
The Input/Output benchmark paradigm restricts the metrics that can be collected during the benchmarking process (see section~\ref{section:32-literature-benchmarking-gai}).

\nrscomment{this is redundant to the previous "This section reviews..." sentence. cut it and incorporate the 3 bottleneck names up in the previous "This section reviews" sentence}
By reviewing workload analyses of \gai systems, and the designs of \gai accelerators we identify that the three driving bottlenecks in \gai systems are \textit{memory, compute} and \textit{communication}.

\nrscomment{cut this sentence}
These are the three bottlenecks that exist in any system, and the following discussion addresses each briefly in turn, before turning to the issue of \textit{accuracy bottlenecks} and \textit{data-dependence}. 

\subsection{Memory Bottlenecks}
Memory bottlenecks occur when a system is unable to keep a processor fully occupied because the data required for computation cannot be moved from memory to the processor quickly enough~\cite{Hennessy2019}.
Graph algorithms are widely understood to be \textit{memory bound}~\cite{Beamer2015}.
Memory bottlenecks manifest either because the volume of data required for processing cannot fit in the memory bandwidth of the system, or because the latency to retrieve data from memory is too long, and the processor spends cycles unoccupied, waiting for data.
Memory bandwidth defines how much data can be retrieved from memory in a single cycle.
When bandwidth is insufficient because of inefficient architectures~\cite{Beamer2015}, poor cache utilization~\cite{Ahn2015}, competing parallel processes~\cite{Fu2014, Buluc2012} or just because the irregular shape of the data leads to them being underutilized~\cite{Aananthakrishnan2020, Zhang2017, Song2016} the processor cannot get enough data to sustain maximal operations and stalls~\cite{Satish2014}.
Memory latency defines how much time it takes to retrieve required data from memory, and exacerbates the under-utilization of memory bandwidth~\cite{Kanev2015, Wu2021a}.
Latency is problematic in \gai applications because the irregular data accesses coming from the inherent lack of spatial and temporal locality (on average)~\cite{Oguntebi2016, Ham2016} causes cache misses~\cite{Song2016, Zhang2017, Beamer2015} and inefficient cache usage even when it is hit~\cite{Ham2016, Mukkara2018, Aananthakrishnan2020}, leading to slow main memory accesses. 
The impact is so significant that even efficient pipelined architectures can't `hide' the latency~\cite{Dysart2016}.

\subsection{Computation Bottlenecks}
Computation bottlenecks occur when the system can serve data to the processor faster than computations can be performed on it. 
Optimizations for general-purpose architectures fail to deliver many benefits because of the irregular data-dependent executions in graph processing~\cite{Dadu2021}. 
For example, branch misprediction is common~\cite{Aananthakrishnan2020}, but is typically eclipsed by memory latency~\cite{Beamer2015} and distributed workloads frequently become unbalanced~\cite{Aananthakrishnan2020, Ozdal2017}.
Looking at GNNs, their hybrid computing patterns that combine irregular graph processing and regular dense neural network processing make it difficult to implement GNNs on general-purpose architectures~\cite{Yan2020}.
The most striking computing problem however is the lack of optimized graph creation, update and delete functionality, with up to 40\% of latency in dynamic graph processing stemming from graph creation~\cite{Basak2020}.

\subsection{Communication Bottlenecks}
Communication bottlenecks occur generally only in parallel computing settings and are caused by separation, movement and synchronization overheads.
Parallel processing requires the data to be divided up between workers. 
Partitioning graphs is a known $NP$-complete problem and the time taken to partition graphs for processing has a significant impact on overall efficiency~\cite{Mccune2015}.
When partitioned, the graph must then be copied to a distributed worker node (or an accelerator) and the time taken to copy the data for processing degrades efficiency~\cite{Mutlu2023, Kiningham2022}.
Synchronizing parallel worker nodes becomes challenging when the algorithm needs to access all vertices each iteration~\cite{Satish2014} or when the partitioning is suboptimal~\cite{Chen2019}.
In either case, distributed parallel approaches incur the most severe communication overheads~\cite{Beamer2015, Dhulipala2021}.

\subsection{Accuracy Bottlenecks}
\nrscomment{this reads like something youre inventing out of thin air. its not really a bottleneck in the sense of being the limiting reagent. I think a better place to put this (and jsut call it the law of diminishing return) is in the benchmark design considerations in ch4}
What we call an accuracy bottleneck occurs when the desired degree of accuracy (or some other effectiveness) metric is beyond the elbow point for diminishing returns, where the rate of improvement levels off proportional to the time and compute applied to it. 
Accuracy bottlenecks are not well explored in \gai benchmarking literature but are a well-known trade-off in machine learning more broadly.
The trade-off is recognized at the macro-level with the recognition that the amount of computing power invested into deep learning is yielding ever smaller returns~\cite{Thompson2021} and at the micro level avoiding overfitting has relevant strategies like early stopping~\cite{Ying2019} that can be employed.
Detecting and acting on accuracy bottlenecks is important if a \gai benchmark is going to be capable of trading off between effectiveness, efficiency and cost. 

\subsection{Data-Dependence}
\nrscomment{this is also kind of arbitrary, and definitely doesnt belong in the bottleneck section. i think you basically covered all of this already in prelim and can just cut it. if you want to keep this as a summary of all the properties that affect efficiency, again move it to design considerations or somewhere in ch4}
There is an inherent dependence between the efficiency and effectiveness of \gai tasks and the size, shape, sparsity, content and velocity of the data they operate on~\cite{Fu2014, Osama2022}.
Large datasets need to be distributed or compressed to fit in memory and the complexity of many \gai tasks requires large graph problems to be parallelized to be tractable~\cite{Hong2012}.
Varying topologies and sparsities can impact a graph's locality properties and memory access patterns~\cite{Lenharth2016, Beamer2015}. 
The content (attributes) of a graph can introduce memory access overheads, but also simplify some algorithms by acting as filters to reduce the size of the graph~\cite{Mccreesh2018,Hamilton2017} and the ever-increasing velocity of data drives the use of dynamic graphs in real-world situations~\cite{Regli2022a, Basak2020, Zang2012}.


% % % % % % % % %  BELOW IS FULL-LENGTH VERSION OF THIS SECTION THAT GOES INTO MUCH MORE DETAIL % % % % % % % %

% A \textit{bottleneck}, or \textit{system constraint} is any resource whose capacity is equal to or less than the demand placed on it~\cite{Goldratt2004}. 
% To improve the overall performance of a system under test, we must identify the bottleneck, exploit that bottleneck to maximize performance with the resources at hand, subordinate the rest of the system to that bottleneck, elevate the capacity of the bottleneck and then re-evaluate the system to find the new bottleneck~\cite{Goldratt1990}.
% This section reviews the prior work in identifying the bottlenecks known to occur during \gai, setting the scene for the discussion of acceleration techniques to \textit{exploit} and \textit{elevate} those constraints in section~\ref{section:34-literature-gai-acceleration}.

% Table~\ref{table:benchmark-metrics} in section~\ref{section:32-literature-benchmarking-gai} shows that most benchmark suites cannot capture the detailed metrics required to identify specific bottlenecks because of limited access to the systems under test to implement code instrumentation and system profiling.
% \textit{Workload Analyses} are the more detailed, invasive investigations that have access to the environment and the implementation to gather the required data. 
% A review of existing workload analyses reveals three key \textit{efficency} bottlenecks exist for \gai applications. 
% \textit{Memory, computation} and \textit{communication} bottlenecks.
% Less explored, but important to consider is the idea of \textit{accuracy tolerances}. 
% An accuracy tolerance determines how accurate a solution needs to be to be acceptable, and is particularly important for approximate or incremental algorithms, where larger computing times may yield more accurate solutions.
% \subsection{Memory Bottlenecks}
% Memory bottlenecks occur when a system is unable to keep a processor fully occupied because the data required for computation cannot be moved from memory to the processor quickly enough~\cite{Hennessy2019}.
% Graph algorithms are widely understood to be \textit{memory bound}~\cite{Beamer2015}. 
% They are \textit{data intensive} algorithms, rather than the \textit{compute intensive} algorithms that modern general-purpose and GPU-accelerated systems are optimized for~\cite{Dysart2016}.
% Memory bottlenecks manifest either because the volume of data required for processing cannot fit in the memory bandwidth of the system, or because the latency to retrieve data from memory is too long, and the processor spends cycles unoccupied, waiting for data.

% \textbf{Memory Bandwidth.} 
% Memory bandwidth defines how much data can be retrieved from memory in a single cycle. 
% In general, being able to load more memory in one cycle (large bandwidth) is better for dense, data-hungry operations.
% Insufficient \textit{memory bandwidth} is commonly identified as the source of memory-boundedness for graph algorithms~\cite{Satish2014} particularly when the instruction window size cannot hold enough instructions to account for cache-missing instructions~\cite{Beamer2015}. \nrscomment{that last sentence doesnt make sense}
% It is a bottleneck because the irregular memory accesses prevent the effective use of the cache hierarchy and because very little computation is performed on each vertex when loaded from memory~\cite{Ahn2015}.
% Na\"{\i}ve `flat' parallel processing implementations exacerbate memory bandwidth issues and impose a scaling limit on the number of processors as they for \nrscomment{youre missing some words here} limited memory bandwidth~\cite{Fu2014}.
% The observation is more concretely characterized by a decay factor of $P^{\frac{1}{2}}$ where $P$ is the number of processors~\cite{Buluc2012} because of the unfavorable communication-computation ratio.
% % While experiments on GPUs have shown that \textit{elevating} the memory bandwidth can increase throughput~\cite{Wu2015} through brute-force alone, it is insufficient to just increase the bandwidth, and more sophisticated methods are needed. 
% However, studies of memory bandwidth utilization in general show that around $95\%$ of production systems utilize less than one-third of their memory bandwidth~\cite{Kanev2015}.
% \gai bottlenecks are no exception with GNNs shown to utilize less than 40\% of their capacity~\cite{Wu2021a} and traditional graph processors inefficently utilizing their memory bandwidth because of inefficent memory reads~\cite{Beamer2015, Aananthakrishnan2020, Zhang2017, Song2016}.
% % with workload analyses showing that there is a plateau in performance even for GPUs with significantly larger memory bandwidths because the sparseness of graph data structures causes inefficient memory reads and wasted bandwidth~\cite{Aananthakrishnan2020, }, especially when large instruction sets further consume limited memory bandwidth~\cite{Beamer2015}.

% \textbf{Memory Latency.} 
% Memory latency defines how much time it takes to retrieve required data from memory. 
% General-purpose architectures alleviate memory latency constraints with the caching system in the memory hierarchy~\cite{Hennessy2019}.
% The cache hierarchy is a collection of small, fast memory modules located close to the processor, optimized to reduce the time it takes to load data from the cache to the processor. 
% The memory hierarchy predicts what data the system will likely require soon because it is proximal in memory to data currently in use (spatially local) or has been used recently and so will likely be needed again (temporally local) and places it in the cache for the processor to use.
% When the data required for the next processing step is not in the cache, the system must wait for it to be retrieved from the main memory, which incurs a much longer delay (latency) 

% Caches are under-utilized in graph processing systems because of irregular memory accesses, resulting in a high cache miss rate~\cite{Song2016} that has been measured as causing $60-80\%$ of cycles to stall waiting for memory accesses~\cite{Zhang2017} and negatively impacting system throughput~\cite{Beamer2015}.
% The ineffectiveness stems from only small portions of each cache line containing the data needed for the next computation step, wasting much of the memory bandwidth that is used to load data into the cache~\cite{Ham2016, Mukkara2018, Aananthakrishnan2020}. 
% The `wasted space' in the cache lines is because of a lack of \textit{spatial locality} in graph data structures, which is generally true for many graph types, but particularly \textit{mesh} graphs with large diameters and a low mean degree~\cite{Oguntebi2016, Ham2016}.
% The lack of locality combined with the data intensity of graph operations results in the latency of memory accesses overwhelming even highly optimized pipelined architectures~\cite{Dysart2016}.

% \subsection{Computation Bottlenecks}
% Computation bottlenecks occur when the system can serve data to the processor faster than computations can be performed on it. 
% While graph processing is typically memory-bound, there are computational bottlenecks that can occur. 

% \textbf{Load Balancing.}
% Load-balancing is a key challenge, particularly when parallelizing graph processing and working with graphs with heavily skewed degree distributions~\cite{Aananthakrishnan2020}.
% The irregular structure of graphs makes it hard to predict what work needs to be completed before execution time.
% For example, many graph algorithms iterate until convergence, however, the number of iterations required to converge individual vertices or subgraphs can vary significantly, blocking other workers while they wait for the over-loaded worker to finish~\cite{Ozdal2017}.

% \textbf{Branching.}
% Graph processing is computationally challenging for traditional architectures~\cite{Dadu2021}.
% In addition to the pre-fetching and cache issues already discussed, graphs can make \textit{branch-prediction} and \textit{out-of-order} processing optimizations ineffective. 
% Graph applications heavily use branches because future processing depends heavily on the structure and order of the input data that cannot be known until input time.
% The branching operations are data dependent, irregular and so are hard to predict which contributes to the underutilization of out-of-order processors~\cite{Aananthakrishnan2020}.
% Luckily, the impact of memory bandwidth under-utilization and memory latency is so significant that the vast majority of the time the branch misprediction rate has no meaningful impact on computation time because it is masked by stalled memory accesses~\cite{Beamer2015}.

% \textbf{Graph Creation.}
% Most \gai evaluations ignore the cost of constructing a graph, considering it a one-off cost and amortizing the cost across the queries that will be conducted on that graph. 
% Ignoring the computation overheads of creating and maintaining graphs is not an option for GDBs and dynamic graphs, which are constantly updating their content.
% Up to $40\%$ of latency in dynamic graph processing comes from the graph update process~\cite{Basak2020}. 
% The impact is emphasized further by efforts to pre-process a graph for optimal execution of PageRank, with $15$ hours of pre-processing to optimize the layout of the graph in memory yielding a reduction in execution time from $12$ seconds to $5$ seconds~\cite{Balaji2018}, requiring almost $7,000$ executions of PageRank to amortize the pre-processing cost. 

% \textbf{Hybrid Computation.}
% GNN-driven components of \gai struggle in particular with hybrid execution patterns. 
% That is, they need to handle the sparse, irregular computations of graph processing as well as the dense, regular computations of Neural Network training and inference~\cite{Yan2020}.
% These hybrid computation patterns can be particularly difficult to accelerate as they require different hardware architectures to fully optimize each component. 

% \subsection{Communications Bottlenecks}
% Communication bottlenecks occur generally only in parallel computing settings.
% Parallel approaches (both distributed and shared-memory) to graph processing are essential to addressing the memory and computation bottlenecks, but they introduce communication bottlenecks because of the poor \textit{Communication to Computation} ratio of graph processing, which is seen as the primary bottleneck impeding the scaling of graphs processing to exa-scale~\cite{Azad2021}.
% The two driving causes of communication bottlenecks are the needs to partition the graph graphs \nrscomment{wording}, copy graph data between worker nodes and the need to enforce synchronization throughout processing.

% \textbf{Graph Partitioning.}
% To parallelize the processing of graph algorithms, protections need to be put in place to prevent different workers from operating on the same data. 
% Traditional parallel approaches will partition the input data, distribute it to workers for processing, and then combine the outputs. 
% In graph processing, an optimal partition of the graph will result in a minimal number of inter-partition dependencies and improve the processing throughput. 
% Optimal partitioning is, however, an $NP$-Complete problem, and so any partitioning must be heuristically done, and will therefore be sub-optimal. 
% The future of practical large-scale distributed graph processing may be related to finding algorithms that process a graph as independent subgraphs, and the performance of that subgraph-centric processing is closely tied to the effectiveness of large-scale graph partitioning, including streaming and distributed partitioning techniques~\cite{Mccune2015}.

% \textbf{Copying graph data.}
% Once a graph is partitioned, the data must be made available to the workers. 
% In distributed systems, this involves copying data between the worker nodes, which depending on the size of the graph, the network bandwidth and network latency can incur significant costs. 
% Shared memory systems do not have the same overhead, but must enforce mechanisms to protect against race conditions. 

% \textbf{Synchronization.}
% Sub-optimal partitioning, irregular computations and irregular memory accesses results in significant synchronization overheads between workers. 
% The problem is particularly pronounced in algorithms that need to process the entire vertex list~\cite{Satish2014} or when the graph is sub-optimally partitioned between workers and when there is high contention for shared vertices~\cite{Chen2019}.
% Synchronization overheads manifest as a large number of main memory accesses in the shared-memory paradigm and the large volume of inter-processor messages in the distributed paradigm~\cite{Beamer2015}.
% In general, message-passing in a distributed parallel system is more expensive than memory accesses in a shared system~\cite{Beamer2015} and given that very large graphs can be processed in shared-memory parallel environments~\cite{Beamer2015, Dhulipala2021}, the conventional wisdom is to avoid distributed processing where possible. 

% \textbf{Architectures.}
% Graph processing workloads have been shown to run poorly on general-purpose parallel processors because the computing patterns of graph processing do not reflect the conventional highly localized workloads they are optimized to support~\cite{Song2016, Kepner2016}.
% Known issues involve incurring overhead when the parallel framework enforces explicit function calls to exchange data between worker nodes~\cite{Dysart2016} and that general-purpose instruction sets like $x86$ bloat the inter-processor traffic with instructions~\cite{Ham2016}. 
% One of the common thrusts in the discussion of parallel graph processing architectures is that a domain-specific architecture is needed to achieve optimal performance. 

% \subsection{Accuracy Thresholds}
% \osullikomment{Something here?}

% \subsection{The impact of datasets}

% The \textit{efficiency} and \textit{effectiveness} of every \gai task depends on the size and shape of the input graph combined with the algorithm being executed~\cite{Fu2014, Osama2022}, and so different algorithmic approaches to the same problem may need to be used depending on if the graph is a \textit{scale-free, uniform-random, mest} or somewhere in-between~\cite{Lenharth2016}.
% The data dependence of graph algorithms is one of the driving reasons that empirical evaluations and workload analyses favor a collection of `real-world' graphs drawn from the domain of intended use.
% % While the shapes of graphs exist on a spectrum, common `classes' of graphs exist. 
% The size, shape, sparsity, content and velocity of a graph all impact the memory, computation, communication and accuracy of a system.
% In the real world, we expect graphs of interest to \gai to be \textit{large, scale-free}, \textit{attributed, heterogeneous} and \textit{dynamic}~\cite{Regli2022a}.

% \textbf{Size.}
% Large graphs can be difficult to store in memory on a single machine. 
% While efficient representations can represent graphs with $3.5$ billion vertices and $128$ billion edges and execute theoretically work-efficient graph algorithms that can be used in search and community detection in a commodity system with a $1$Tb of Random Access Memory (RAM)~\cite{Dhulipala2018}, estimates of uncompressed commercial graph sizes exceed 20TB~\cite{Aananthakrishnan2020} and the largest Graph500 benchmark comes in at over $60$ terabytes, which is well beyond the capacity of any commercial system. 
% So, for larger graphs and more complicated algorithms, single-machine solutions may not be sufficient, which may drive the decision to employ a distributed computing model, introducing communication bottlenecks that may not otherwise exist. 
% Many \gai tasks have complexity scales that are linear, quadratic or worse in the number of edges and so scale poorly as inputs reach larger sizes~\cite{Hong2012}.
% Therefore, larger, denser graphs tend to mean longer computation times. 

% \textbf{Shape.}
% Graphs with regular degree distributions and large diameters have more regular processing patterns and are more easily partitioned but have poor spatial locality.
% Scale-free graphs, on the other hand, have been shown to improve temporal locality in graph processing, but are more difficult for parallel algorithms to process because they are difficult to cleanly partition and the variance in neighborhood size leads to imbalanced workload between worker nodes~\cite{Beamer2015}.

% \textbf{Sparsity.}
% % The discussion of size, degree, diameter, transitivity and community structure highlights that it is the \textit{topology} of a graph that contains much of the latent information that graph analysis algorithms try to extract. 
% A defining feature of graphs is their \textit{sparsity}.
% \textit{Sparsity} is most easily understood by looking at an adjacency matrix (figure~\ref{figure:adjacency_matrix}).
% Of the $|V|^{2}$ matrix, only $|E|$ of the values contain `information', with the remainder of the values being $0$.
% There are two impacts relevant to \gai. 
% First, if the graph is stored as an adjacency matrix, it also needs to allocate memory to represent the $0$s, greatly increasing the memory footprint. 
% Second, if the graph is stored in the canonical format, even if the $0$s are not represented, many memory management approaches will create the graphs in memory in the order that the vertexes are created, potentially leading to vertices connected by edges (neighbors) being stored in different locations in memory, and so will not be pre-fetched together when being processed. 
% More sophisticated graph database systems use the \textit{index-free adjacency principle} to colocate adjacent nodes in memory and use internal pointers to track the locations of neighbors without a separate index~\cite{Vukotic2015}.

% \textbf{Content}
% Many real-world graphs we are interested in processing are knowledge graphs that are attributed, directed, and weighted and be multigraphs with hyperedges. 
% % A knowledge graph is typically built on a constrained ontology that defines the entities and relations at an abstract level. 
% The information contained in vertex and edge labels is critical to answering search and reachability queries (what is the `name' of the target vertex), can reduce complexity in subgraph isomorphism~\cite{Mccreesh2018}, improve the \textit{effectiveness} of community detection~\cite{Zhang2020} and is critical to learning useful embeddings for GNNs~\cite{Hamilton2017}. 
% However, beyond just increasing the size of the graph, the label information is typically held separately in memory from the graph topology which can impact retrieval times, and the checking and comparison of complicated node attributes can increase the computational intensity of operations. 

% \textbf{Velocity.}
% Much of traditional \gai work assumes a static graph that amortizes the cost of construction across the queries run over its lifetime.
% Dynamic graphs reflect most real-world applications and are in a constant state of change. 
% While the use of static graph processing algorithms in a dynamic context is treated suspiciously~\cite{Basak2020}, they are often all that is available.
% Graphs are updated in batches and performance is measured in batch latency, the time it takes from when a vertex is pushed to the graph until it appears in the result of a \gai task~\cite{Zhang2012}.


% \subsection{Summary of \gai bottlenecks}
% \osullikomment{Summary Here}


% % % % % % % % %  ABOVE IS FULL-LENGTH VERSION OF THIS SECTION THAT GOES INTO MUCH MORE DETAIL % % % % % % % %


%%%%%%%%% OLDER MATERIAL BELOW HERE %%%%%%%%%%

% The sizes of graphs have been increasing in the `big data' era, w
% Beyond the Input/Output (I/O) problems that `big data' graphs pose in the face of no universally accepted binary exchange format~\cite{Azad2021}, large graphs quickly became unable to fit in the memory of a single machine, and some algorithms 
% When graphs are unable to fit in memory, the choices become \textit{distributing} the computation, or \textit{compressing} the representation~\cite{Shun2015, Mccune2015}. 
% Popular compressed representations like the CSR/CSC formats will become hyperspace when graph decomposition techniques are employed~\cite{Azad2021}, so more aggressive compression is required to make it a useful tactic.
% Unfortunately, compression is not viable for sequential algorithms, as the encoding/decoding cost for the compressed representation tends to outweigh the benefit of the smaller representation~\cite{Shun2015}, forcing the use of parallel implementations to `hide' the cost.
% Shared memory parallelism has proven able to represent and process graphs with over $3.5$ billion vertices and $128$ billion edges within $1$ TB of RAM~\cite{Dhulipala2021}, demonstrating its viability as an alternative to distributed parallelism. 



% Distributed computing has proved effective for other large-scale problems, and so is a natural path to follow for large-scale graph analysis. 
% Unfortunately, graph algorithms perform poorly in a distributed setting. 
% There are two main reasons. 
% First, divining the work optimally involves partitioning the graph optimally. 
% Graph partitioning is itself an $NP$-complete problem, and so heuristic partitioning requires some up-front computation, and may still yield a sub-optimal result. 
% \osullikomment{Cite Graph Partition NP Complete}
% Second, because of the typically \textit{scale-free} characteristic of I\S graphs, we see that for the densely connected parts of the graph, even an optimal partitioning may have cross-partition edges. 
% Consequently, the control mechanisms used to synchronize the processing may need to `wait' to access information on the other machine. 
% It also introduces delay as network bandwidth and round-trip latency degrade the communication channel further. 
% \osullikomment{cite distributes dramas.}





% The majority of datasets of interest to the I\&S community share four features. 








% Overall, the processing of a graph is highly dependent on the shape of its input data. 
% \textit{Scale-free} graphs have better temporal locality than \textit{mesh} graphs, but both struggle with optimizing for spatial locality~\cite{Beamer2015}. 
% Combined, the lack of spatial and temporal locality results in poor utilization of the memory hierarchy in the average case~\cite{Balaji2018}.
% % The ineffective use of the memory hierarchy is exacerbated by the increasing sizes of graphs, with the largest not fitting in the memory of a single machine, and a known negative correlation between graph size and cache utilization. 
% Graph processing is memory-bound and finding ways to \textit{Exploit} and \textit{Elevate} remain important research priorities. 

% \textbf{Graph Size and Shape.} 




% \textit{Workload Analyses} are detailed reports and empirical studies of how \textit{systems} execute various \textit{workloads}. 
% They are typically more involved, involving code instrumentation and system profiling to build a detailed understanding of computational, memory and network bottlenecks. 




% \nrscomment{idk if we have the background to really grasp this point yet. maybe bring it up when you talk about representation and spatial locality? at this stage i dont know why scale free would improve locality, since i dont even know how graphs are stored yet}

% Mesh Graphs are typically better suited to parallel processing because their large diameter and local neighborhood structure mean there are minimal cross-partition dependencies, and the workload is more balanced. 
% However, for algorithms that must visit the entire graph each iteration (like PageRank), the poor spatial locality of the vertices in memory incurs additional processing latency. 
% The expressive power of graphs is being explored beyond `graph-native' problems. 
% \nrscomment{same issue here, dont really have the background to talk about this yet}



% \nrscomment{this characteristics section feels like it should be moved to the abstract graph benchmark design. these are basically your characteristics to vary and account for. it doesnt belong up here. rest of the comment sin this section still apply but assume that this whole chunk will be with the benchmark stuff later.}

% \subsection{General Characteristics}
% \label{section:124-general-characteristics}

% Decades of research yield the characteristics of graphs that are `generally true' and interact with the \textit{efficiency} of graph analytics.
% These graph characteristics focus on the \textit{topological} features of graphs.
% Depending on the context, the addition of attribute information may simplify or complicate the algorithms that operate on the graphs. 

% \textbf{Input dependence.} 
% %
% \nrscomment{"Input dependence" is not a graph characteristic, and this subsubsection talks about random graphs. I think you can just cut this chunk entirely, or maybe reframe it, since you keep bringing up random graphs when explaining the subsequent characteristics. Maybe you first describe broad classes of graph domains (synthetic vs real. then within that the hierarchy you drew out) and compare the 'naive' random graph that one might use for a benchmark to the others and highlight the characteristics that matter for the benchmark. For instance show random edges vs fully connected, visually, and then explain why it affects processing. And then repeat for the other characteristics.}
% %





% \subsubsection{Summary of Graph Characteristics}
% \osullikomment{Summary here}


% \subsubsection{Dataset Characteristics}


% % 



% \subsubsection{Communication Bottlenecks}


% \subsubsection{Computation Bottlenecks}




    



% \textbf{Memory Boundedness.}
% Graph algorithms are widely understood to be \textit{memory bound}~\cite{Beamer2015}. 
% They are \textit{Data Intensive} algorithms, rather than the \textit{Compute Intensive} algorithms that modern general-purpose and GPU-accelerated systems are optimized for~\cite{Dysart2016}.
% A program is \textit{memory bound} when the processor is not fully utilized, typically because it is stalled waiting for data to be retrieved from memory.
% The \textit{memory bandwidth} of a system is commonly identified as the source of memory-boundedness for both \textit{traversal} and \textit{computationally intensive} graph algorithms~\cite{Satish2014}.
% Memory bandwidth is a pressure point because the large number of irregular memory accesses prevent the effective use of the cache hierarchy and because very little computation is performed on each vertex when loaded from memory~\cite{Ahn2015}.

% \textbf{The Memory Hierarchy.} General-purpose architectures alleviate memory bandwidth constraints with the caching system in the memory hierarchy~\cite{Hennessy2019}.
% Caches are under-utilized in graph processing systems because of irregular memory accesses, resulting in a high cache miss rate~\cite{Song2016} that has been measured as causing $60-80\%$ of cycles to stall waiting for memory accesses~\cite{Zhang2017}.
% The ineffectiveness of pre-fetching data in graph processing puts additional pressure on the cache hierarchy.
% Only small portions of each cache line contain the data needed for the next computation step, wasting much of the memory bandwidth that is used to load data into the cache~\cite{Ham2016, Mukkara2018, Aananthakrishnan2020}. 
% The `wasted space' in the cache lines is because of a lack of \textit{spatial locality} in graph data structures, which is generally true for many graph types, but particularly \textit{mesh} graphs with large diameters and a low mean degree~\cite{Oguntebi2016, Ham2016}.

% %
% \nrscomment{ok, in that case you need a section explicitly outlining those classes, and possibly 
% a table comparing them across the characteristics you're about to describe. you currently have it mixed in with the characteristics, so i come away not understanding what the classes are or what the characteristics are. decouple them.}
% %

% \textbf{Size Matters.}
% %
% \nrscomment{please a different header}
% \nrscomment{replace this entire description with a picture. this can be 2 sentences}
% %
% Section~\ref{section:11-Intelligence-Security-Context} identifies that graphs in the I\&S domain can be billions to trillions of vertices and edges. 
% The \textit{size} of a graph impacts graph processing. 
% In the most obvious sense, like all computations, there is some relationship between the size of the input and the time required to process it.
% The count of the vertices can give us a false understanding of the scale of the graph. 
% For example, suppose a graph with $10$ vertices. 
% Each vertex except for the first and last is connected by an undirected edge to only two other vertices, such that they form a straight line of $10$ vertices.
% The graph will have $10$ vertices and $9$ edges. 
% But, suppose a fully connected graph, where each of those vertices is connected to every other vertex.
% That then becomes a graph of $10$ vertices and $45$ edges.
% If directed (and symmetric), we would have $18$ and $90$ edges respectively for the same graph of $10$ vertices. 
% If it is a more complicated multi-graph structure the edge count could increase even further. 
% So, any discussion of graph `size' must also account for the number of edges. 
% \osullikomment{Needs a Cite}

% \textbf{Degree Distribution.}
% The \textit{size} of the graph is closely linked to the number of edges in it. 
% Beyond a raw edge count, consideration must be given to the Degree $K$, the mean degree $\bar{K}$, the mean in-degree $\bar{K}_{in}$ and mean out-degree $\bar{K}_{out}$ and the \textit{distribution} of these values across the graph. 
% %
% \nrscomment{we dont know what any of those things mean, or why 'consideration must be given' to them}
% %
% Random graphs are quite regular in their distribution of degrees, and we expect the modal degree to approximate the mean degree.
% Similarly, on graphs derived from physical phenomena like road networks, the mean degree is generally consistent.
% Conversely, graphs derived from social interactions like phone call records or follower lists on social media tend to follow a power-law distribution. 
% %
% \nrscomment{if youre going to define scale-free, do it properly}
% %
% In these \textit {`Scale-Free'} graphs only a small percentage of the vertices possess most of the edges, and most vertices have very few edges.
% Intuitively, consider that a few celebrities are extremely popular and have large followings, while most people have only small followings.
% The scale-free graph shape is one of the key drivers of computational \textit{inefficeny} in graph processing. (see: section~\ref{section:126-processing-bottlenecks}).
% \osullikomment{Needs a Cite}

% \textbf{Graph Diameter and Transitivity}
% The \textit{diameter} of the graph is the length of the shortest path between the two most `distant' nodes.
% The \textit{transitivity} describes the probability that adjacent nodes in a graph are connected. 
% The \textit{clustering coefficient} is a measure of the transitivity at a local or global level (with slightly different calculations depending on scope).

% Random graphs tend to have a middling diameter and clustering coefficient, depending on the density parameter used to generate the graph.
% Graphs derived from physical phenomena like road networks tend to have a large diameter and small global clustering coefficient.
% Scale-free graphs representing social phenomena like social media follower graphs tend to have a small diameter and large clustering coefficient. 
% Known as the \textit{small-world} characteristic, we can say that of a scale-free graph, most vertices are `close' to each other.
% Some graphs have a \textit{small-world} property in some component subgraphs but have other subgraphs that are of a larger diameter and lower density.

% The variance in the diameter and density of graphs, even internally on large complicated graphs like \textit{knowledge graphs} means that a graph processing algorithm may need to intermittently deal with extremely dense and sparse sections within the same algorithm, complicating the ability to efficiently process them.
% \osullikomment{Needs a Cite}


% \textbf{Community Structure.}
% Closely linked to the idea of degree distribution and clustering coefficient is \textit{community structure}.
% Characterized by Michelle Girvan Mark and Newman in 2002, a community is broadly defined as a subgraph that is more densely connected to itself than it is to the remainder of the parent graph~\cite{Girvan2002}.
% It is a generalization of \textit{Cliques}, which are fully connected subgraphs within a parent graph.
% Communities are more distinct in graphs that have variance in their degree distributions and local versus global structure. 
% Exploiting the latent information in communities relies on the variance in graph structure and provides a useful highlight of why graphs of heterogeneous shape are not only expected but essential in many of the graph analysis tasks. 


% A \textit{bottleneck}, or \textit{system constraint} is any resource whose capacity is equal to or less than the demand placed on it~\cite{Goldratt2004}. 
% To improve the overall performance of a system, we must identify the bottleneck, exploit that bottleneck to maximize performance with the resources at hand, subordinate the rest of the system to that bottleneck, improve the bottleneck and then re-evaluate the system to find the new bottleneck~\cite{Goldratt1990}.
% This section reviews the prior work in identifying the bottlenecks known to occur during graph processing, setting the scene for the discussion of using hardware accelerators to \textit{exploit} and \textit{elevate} those constraints in section~\ref{section:13-Graph-Hardware-Acceleration}.


% \subsubsection{Memory Bottlenecks}
% \textbf{Graph Size and Shape.} 
% The sizes of graphs have been increasing in the `big data' era, with estimates of uncompressed commercial graph sizes exceeding 20TB~\cite{Aananthakrishnan2020}.
% Beyond the Input/Output (I/O) problems that `big data' graphs pose in the face of no universally accepted binary exchange format~\cite{Azad2021}, large graphs quickly became unable to fit in the memory of a single machine, and some algorithms scale poorly as graphs reach larger sizes~\cite{Hong2012}.
% When graphs are unable to fit in memory, the choices become \textit{distributing} the computation, or \textit{compressing} the representation~\cite{Shun2015, Mccune2015}. 
% Popular compressed representations like the CSR/CSC formats will become hyperspace when graph decomposition techniques are employed~\cite{Azad2021}, so more aggressive compression is required to make it a useful tactic.
% Unfortunately, compression is not viable for sequential algorithms, as the encoding/decoding cost for the compressed representation tends to outweigh the benefit of the smaller representation~\cite{Shun2015}, forcing the use of parallel implementations to `hide' the cost.
% Shared memory parallelism has proven able to represent and process graphs with over $3.5$ billion vertices and $128$ billion edges within $1$ TB of RAM~\cite{Dhulipala2021}, demonstrating its viability as an alternative to distributed parallelism. 
% Together, the size and shape of the input graph, combined with the algorithm being executed has a significant impact on computing efficency~\cite{Fu2014, Osama2022}, and so different algorithmic approaches to the same problem may need to be used depending on if the graph is \textit{random, scale-free} or \textit{mesh}~\cite{Lenharth2016}.

% \textbf{Memory Boundedness.}
% Graph algorithms are widely understood to be \textit{memory bound}~\cite{Beamer2015}. 
% They are \textit{Data Intensive} algorithms, rather than the \textit{Compute Intensive} algorithms that modern general-purpose and GPU-accelerated systems are optimized for~\cite{Dysart2016}.
% A program is \textit{memory bound} when the processor is not fully utilized, typically because it is stalled waiting for data to be retrieved from memory.
% The \textit{memory bandwidth} of a system is commonly identified as the source of memory-boundedness for both \textit{traversal} and \textit{computationally intensive} graph algorithms~\cite{Satish2014}.
% Memory bandwidth is a pressure point because the large number of irregular memory accesses prevent the effective use of the cache hierarchy and because very little computation is performed on each vertex when loaded from memory~\cite{Ahn2015}.

% \textbf{The Memory Hierarchy.} General-purpose architectures alleviate memory bandwidth constraints with the caching system in the memory hierarchy~\cite{Hennessy2019}.
% Caches are under-utilized in graph processing systems because of irregular memory accesses, resulting in a high cache miss rate~\cite{Song2016} that has been measured as causing $60-80\%$ of cycles to stall waiting for memory accesses~\cite{Zhang2017}.
% The ineffectiveness of pre-fetching data in graph processing puts additional pressure on the cache hierarchy.
% Only small portions of each cache line contain the data needed for the next computation step, wasting much of the memory bandwidth that is used to load data into the cache~\cite{Ham2016, Mukkara2018, Aananthakrishnan2020}. 
% The `wasted space' in the cache lines is because of a lack of \textit{spatial locality} in graph data structures, which is generally true for many graph types, but particularly \textit{mesh} graphs with large diameters and a low mean degree~\cite{Oguntebi2016, Ham2016}.


% \textbf{Memory Bandwidth Utilization.} 
% Na\"{\i}ve `flat' parallel processing implementations exacerbate memory bandwidth issues and impose a scaling limit on the number of processors as they for limited memory bandwidth~\cite{Fu2014}.
% The observation is more concretely characterized by a decay factor of $P^{\frac{1}{2}}$ where $P$ is the number of processors~\cite{Buluc2012} because of the unfavorable communication-computation ratio.
% However, studies of memory bandwidth utilization in general show that around $95\%$ of production systems utilize less than one-third of their memory bandwidth~\cite{Kanev2015}.
% Graph-specific studies indicate that the existing bandwidth is poorly utilized when the instruction window size cannot hold enough instructions to account for cache-missing instructions~\cite{Beamer2015}.
% While experiments on GPUs have shown that \textit{elevating} the memory bandwidth can increase throughput~\cite{Wu2015} through brute-force alone, it is insufficient to just increase the bandwidth, and more sophisticated methods are needed. 

% \textbf{Memory Latency \& Attributed Graphs.} 
% Contributing to the memory bandwidth issue, memory latency has a significant negative impact on system throughput~\cite{Beamer2015}.
% The lack of locality combined with the data intensity of graph operations results in the latency of memory accesses overwhelming even highly optimized pipelined architectures~\cite{Dysart2016}.
% The latency problem is exacerbated by attributed graphs. 
% It is common for the topological and attribute data of a graph to be stored separately in memory~\cite{Buluc2013}. 
% Figure~\ref{figure:locality-attributed} expands on the example from figure~\ref{figure:locality-simple} by adding in vertex attributes. 
% The attributes are stored at separate memory locations, with pointers from the topological vertex (memory positions $1-9$) to their respective attributes (memory positions $10-29$). 
% Attributes consume more memory so more memory bandwidth is required to load them.
% The heterogeneous data types being loaded in attributed graphs (like knowledge graphs) have been shown to inflict significant Dynamic Random Access Memory (DRAM) accesses and hence, increase latency and memory bandwidth pressure~\cite{Basak2019}.
% To illustrate why, again consider a BFS that is searching for the `\textbf{Grape}' vertex. 
% Each time a vertex is dequeued, its attribute must be retrieved from memory to see if it is the target vertex. 
% In the worst case, this is a memory access each time a vertex is checked, or, if reducing the latency of the comparison by pre-fetching the attribute along with the vertex, the cache block of $3$ quickly becomes filled with the vertex plus its attribute, exacerbating the spatial locality problem.

% Overall, the processing of a graph is highly dependent on the shape of its input data. 
% \textit{Scale-free} graphs have better temporal locality than \textit{mesh} graphs, but both struggle with optimizing for spatial locality. 
% Combined, the lack of spatial and temporal locality results in poor utilization of the memory hierarchy in the average case~\cite{Balaji2018}.
% The ineffective use of the memory hierarchy is exacerbated by the increasing sizes of graphs, with the largest not fitting in the memory of a single machine, and a known negative correlation between graph size and cache utilization. 
% Graph processing is memory-bound and finding ways to \textit{Exploit} and \textit{Elevate} remain important research priorities. 


% \subsection{Processing Bottlenecks}
% \label{section:126-processing-bottlenecks}

% \osullikomment{Intent is to use the exploit and elevate framwework i introduce here to taxonomize the acceleration methods, probably generate a table or taxonomic graph of them for a pretty visual.}

% Graph processing began, as most computer science problems do, as a math problem--the K{\"o}nigsberg bridge problem--asking if the seven bridges of K{\"o}nigsberg could be traversed in a single trip without doubling back, ending back at the start point. 
% Euler's answer to this problem is generally heralded as the birth of graph theory~\cite{Barnett2005}.
% The early computational approaches to graph problems focused on sequential algorithms on single machines. 
% The rise of the `big-data' era has seen graphs increase in size from the four vertexes and seven edges of the K{\"o}nigsberg bridge problem to billions or trillions of nodes and edges. 

% The increased graph size has flow-on effects. 
% Large graphs quickly became unable to fit in the memory of a single machine, and even where they could because of a powerful machine or a compressed graph representation, sequential algorithms were too slow to process the graph in a reasonable time. 
% Distributed computing offered a promising path to \textit{elevate} the single machine capacity bottleneck, but problems of how to efficiently partition graphs and synchronize operations on those partitions for distributed processing imposed new constraints, with network bandwidth, network latency and stalled cycles waiting for synchronization degrading performance. 

% As computing capacity increased, \textit{elevating} the memory bottleneck, shared memory parallelism for graph processing gained traction.
% With viable shared-memory parallel solutions emerging, the focus of the research switched to improving the efficiency of programming frameworks and the processing itself to \textit{exploit} these gains. 
% The large number of graph-processing frameworks that exist `to abstract the complexity of implementing complicated parallel graph applications' evidences the growing demand for graph processing. 
% Realizing that the performance of most graph processing algorithms is dependent on the shape of the input data drove a wave of efforts to \textit{exploit} the interaction of the algorithm with the representation by re-ordering the graph as a pre-processing step, scheduling the order that vertices are ordered, and partitioning the graph so subgraphs can fit into cache. 

% Deeper analysis shows that graph algorithms are almost universally memory-bound, except for some distributed implementations which are network-bound. 
% The constraints are commonly attributed to memory bandwidth and memory latency emerging from poor exploitation of the cache hierarchy and pre-fetching optimizations by graphs on CPU-based single machines, and a high communication-computation ratio for parallel (distributed and shared-memory) implementations. 
% Hardware acceleration is a natural step to \textit{elevate} both local and distributed approaches. 
% Representing graphs as matrices can allow us to \textit{exploit} the significant memory bandwidth available on GPUs, though sparsity quickly becomes a problem for the efficiency of both representation and computation, as well as incurring overhead for copying the graph data from main memory to the much more limited on-chip GPU memory. 
% Other approaches use FPGAs and ASICs to \textit{elevate} the computation bottlenecks in graph processing. 
% However, to maximize efficiency, they need to be hand-tuned to both the distribution of the input graph and the algorithm being deployed on the system.

% There is a particular emphasis on accelerating GNN hardware as a special case of graph processing made more difficult by the \textit{hybrid} execution patterns that include both the sparse, irregular memory accesses of graph processing and the dense, regular computation of neural networks.
% Taken together we can see that the efficient \textit{general purpose} processing of graph data structures is an open problem despite decades of work to \textit{exploit} graph processing with algorithmic and data structure innovations and \textit{elevate} performance through parallelization and hardware acceleration. 
% Estimates of commercially interesting graph sizes range up to $20$ Terabytes (TB)~\cite{Aananthakrishnan2020} and as datasets become larger, approaching the exa-scale, we are likely to see an increase in the demand for accelerated, distributed graph processing~\cite{Azad2021}.

% The following sections summarize each of these advances in turn and highlight common themes in the limitations that are used to inform the benchmark framework presented in chapter~\ref{chapter:30-Graph-Hardware-Benchmarking}.
